% Part: lambda-calculus
% Chapter: syntax
% Section: free-variables

\documentclass[../../../include/open-logic-section]{subfiles}

\begin{document}

\olfileid[fr]{lam}{syn}{fv}
\olsection{Variables libres}

Le lambda-calcul porte sur les fonctions,
et l'abstraction lambda permet de les former.
Intuitivement, $\lambd[x][M]$ est la fonction
dont les valeurs sont données par~$M$ lorsqu'on
attribue l'argument de cette fonction à~$x$.
Mais toutes les occurrences de~$x$ dans~$M$
ne jouent pas ce rôle : si $M$ contient une autre
abstraction $\lambd[x][N]$, les occurrences de $x$
dans~$N$ relèvent de $\lambd[x][N]$, et non de
$\lambd[x][M]$. Ainsi, le $\lambd[x]$ initial de
$\lambd[x][M]$ \emph{lie} les occurrences de $x$
dans~$M$ que ne lie pas déjà une autre
abstraction : les occurrences \emph{libres}
de~$x$ dans~$M$.

\begin{defn}[Portée]
Si $\lambd[x][M]$ figure dans un terme~$N$,
l'occurrence correspondante de~$M$ est la
\emph{portée} du~$\lambd[x]$.
\end{defn}

Par exemple, dans $(\lambd[x][x])y$, la portée de
$\lambd[x]$ se réduit à l'occurrence de $x$ dans
son corps. L'argument $y$ est en dehors de cette portée.

\begin{defn}[Occurrence libre et occurrence liée]
Une occurrence de la variable $x$ (comme sous-terme, et non comme paramètre d'une abstraction) dans un
terme~$M$ est \emph{libre} si elle ne se trouve
dans la portée d'aucun $\lambd[x]$, et
\emph{liée} dans le cas contraire. Une occurrence
de la variable~$x$ située dans le corps de
$\lambd[x][M]$ est liée par le $\lambd[x]$
initial si et seulement si cette occurrence de
$x$ est libre dans~$M$.
\end{defn}

\begin{ex}
Dans $\lambd[x][x y]$, $x$ et $y$ sont tous
deux dans la portée de $\lambd[x]$; $x$ est
donc lié par $\lambd[x]$. Comme $y$ n'est
dans la portée d'aucun $\lambd[y]$, il est
libre. Dans $\lambd[x][x x]$, les deux
occurrences de $x$ sont liées par~$\lambd[x]$,
car toutes deux sont libres dans $xx$. Dans
$((\lambd[x][xx])x)$, la dernière occurrence
de~$x$ est libre, car elle n'est dans la portée
d'aucun $\lambd[x]$. Dans
$\lambd[x][(\lambd[x][x])x]$, la portée du
premier $\lambd[x]$ est $(\lambd[x][x])x$,
et celle du second $\lambd[x]$ est
l'avant-dernière occurrence de~$x$. Dans
$(\lambd[x][x])x$, la dernière occurrence de
$x$ est libre et l'avant-dernière est liée.
Ainsi, l'avant-dernière occurrence de $x$
dans $\lambd[x][(\lambd[x][x])x]$ est liée
par le second $\lambd[x]$, et la dernière,
par le premier $\lambd[x]$.
\end{ex}

Pour un terme $P$, on peut examiner toutes
ses occurrences de variables et former
l'ensemble de ses variables libres. Cet
ensemble se note $\FV{P}$ et se définit
naturellement comme suit.

\begin{defn}[Variables libres d'un terme] \ollabel{def:fv}
L'ensemble des \emph{variables libres} d'un
terme est défini inductivement par :
\begin{enumerate}
  \item $\FV{x} = \{x\}$ \ollabel{def:fv1}
  \item $\FV{\lambd[x][N]} = \FV{N} \setminus \{x\}$ \ollabel{def:fv2}
  \item $\FV{PQ} = \FV{P} \cup \FV{Q}$ \ollabel{def:fv3}
\end{enumerate}
\end{defn}

\begin{prob}
\begin{enumerate}
\item Déterminer les portées de $\lambd[g]$
  et des deux $\lambd[x]$ dans le terme
  \[
    \lambd[g][(\lambd[x][g (x x)]) \lambd[x][g (x x)]]\,.
  \]
\item Dans $\lambd[g][(\lambd[x][g (x x)]) \lambd[x][g (x x)]]$,
  toutes les occurrences de variables sont-elles
  liées? Par quelles abstractions le sont-elles
  respectivement?
\item Déterminer
  $\FV{\lambd[x][(\lambd[y][(\lambd[z][x y]) z]) y]}$.
\end{enumerate}
\end{prob}

\begin{explain}
Une variable libre est comme une référence au
monde extérieur (l'\emph{environnement}). Un
terme contenant des variables libres peut être
considéré comme partiellement spécifié, car son
comportement dépend de la façon dont cet
environnement est fixé. Par exemple, dans le
terme $\lambd[x][f x]$, qui reçoit un
argument~$x$ et renvoie le résultat de
l'application de~$f$ à cet argument, la
variable~$f$ est libre. La valeur de ce terme
dépend donc de son environnement, en particulier
de la valeur de $f$ dans cet environnement.

Si l'on abstrait ce terme une fois de plus,
on obtient $\lambd[f][\lambd[x][f x]]$.
Ce terme ne dépend plus de la variable
d'environnement $f$, car il désigne désormais
une fonction qui reçoit successivement deux
arguments et renvoie le résultat de
l'application du premier au second. Modifier
$f$ dans l'environnement n'aura aucun effet
sur le comportement de ce terme : celui-ci
utilise seulement la valeur reçue comme argument,
et non la valeur de $f$ dans l'environnement.
\end{explain}

\begin{defn}[Terme clos, combinateur]
Un terme sans variable libre est appelé
\emph{terme clos} ou \emph{combinateur}.
\end{defn}

\begin{lem}\ollabel{lem:fv}
\begin{enumerate}
  \item \ollabel{lem:fv-abs} Si $y \neq x$, alors
    $y \in \FV{\lambd[x][N]}$ si et seulement si
    $y \in \FV{N}$.
  \item \ollabel{lem:fv-app} $y \in \FV{PQ}$ si et
    seulement si $y \in \FV{P}$ ou
    $y \in \FV{Q}$.
\end{enumerate}
\end{lem}

\begin{proof}
Exercice.
\end{proof}

\begin{prob}
Démontrer le \olref[lam][syn][fv]{lem:fv}.
\end{prob}

\end{document}
